\documentclass{beamer}

%%%%%%%%%%%%%%%%%%%%%%%%%%%%%%%%%%%%%%%%%%%%%%%%%%%%%%%%%%%%%%%%%%%%%%%%%%%%%%%%
% Packages
%%%%%%%%%%%%%%%%%%%%%%%%%%%%%%%%%%%%%%%%%%%%%%%%%%%%%%%%%%%%%%%%%%%%%%%%%%%%%%%%

% Warn about commands, classes and packages that are outdated and superseded.
% Orthodox checks for pitfalls that are not technically incorrect.
%\usepackage[l2tabu,orthodox,abort]{nag}
\usepackage[l2tabu,orthodox]{nag}

% For font and input encoding.
\usepackage[utf8]{inputenc}
%\usepackage[T1]{fontenc}

% Pretty much all of the ams maths packages.
\usepackage{amsmath,amsthm,amssymb,amsfonts}

% Allows inclusion of graphics easily and "configurably".
\usepackage{graphicx}

% Typesets URLs sensibly --- with a TrueType font, hyperlinked in PDFs, and not
% breaking across lines.
\usepackage{url}

% Count the number of used registers (counter, dimen, skip, muskip, box, token,
% input, output, math families, languages, insertions). Also time the amount of
% time required for LaTeX compilation.
\usepackage[timer,left]{regstats}

% The listings package is a source code printer for LaTeX. You can typeset stand
% alone files as well as listings with an environment similar to verbatim as
% well as you can print code snippets using a command similar to \verb.
\usepackage{listings}

% This package implements high-level interface commands for class file writers
% which allows the separation of formatting commands (typically instances of
% so-called `templates') and their arguments from the signature of
% document-level commands.
\usepackage{xparse}

% This package allows verbatim text to be typeset in bold.
\usepackage{alltt}
\renewcommand{\ttdefault}{txtt}

% This style option improves the interface for defining floating objects such
% as figures and tables in LaTeX. It adds the notion of a `float style' that
% governs appearance of floats. New kinds of floats may be defined using a
% \newfloat command analogous to \newtheorem. This style option also
% incorporates the functionality of David Carlisle's style option here, giving
% floating environments a [H] option which means `PUT IT HERE' (as opposed to
% the standard [h] option which means `You may put it here if you like').
\usepackage{float}

% This is the style file for the gnuplot PGF/TikZ terminal. It is associated
% with the 'gnuplot.lua' script.
\usepackage{gnuplot-lua-tikz}

% The caption package provides many ways to customise the captions in floating
% environments such figure and table.
\usepackage{caption}

% This package supports typesetting of sub-captions (by using the the
% sub-caption feature of the caption package).
\usepackage{subcaption}

% This package reimplements the LATEX verbatim and verbatim* environments. In
% addition it provides a comment environment that skips any commands or text
% between \begin{comment} and the next \end{comment}. It also defines the
% command verbatim to input a whole file verbatim.
\usepackage{verbatim}

% TODO
\usepackage[noline,ruled]{algorithm2e}

% TODO
\usepackage{pgfplots}
\usepackage{tikz}\usetikzlibrary{positioning,shapes,shadows,arrows}

% For \ifthenelse command
\usepackage{ifthen}

% For \maxsizebox command
\usepackage{adjustbox}

% My custom packages
\usepackage{escape,persons,pgf-pie}

%%%%%%%%%%%%%%%%%%%%%%%%%%%%%%%%%%%%%%%%%%%%%%%%%%%%%%%%%%%%%%%%%%%%%%%%%%%%%%%%
% Beamer configuration.
%%%%%%%%%%%%%%%%%%%%%%%%%%%%%%%%%%%%%%%%%%%%%%%%%%%%%%%%%%%%%%%%%%%%%%%%%%%%%%%%
% Presentation mode
\mode<presentation>

% Load a theme
\usetheme{Warsaw}
%\useoutertheme[footline=authortitle]{miniframes}
\useoutertheme{smoothbars}

% For animations in beamer
\usepackage{multimedia}

% Reduce the font size of captions
\setbeamerfont{caption}{size=\scriptsize}

% Show notes
\setbeameroption{show notes}
%\setbeameroption{show only notes}
%\setbeameroption{hide notes}
%\usepackage{pgfpages}\setbeameroption{show notes on second screen}

% Hide navigation symbols
%\setbeamertemplate{navigation symbols}{}

%%%%%%%%%%%%%%%%%%%%%%%%%%%%%%%%%%%%%%%%%%%%%%%%%%%%%%%%%%%%%%%%%%%%%%%%%%%%%%%%
% Custom colours
%%%%%%%%%%%%%%%%%%%%%%%%%%%%%%%%%%%%%%%%%%%%%%%%%%%%%%%%%%%%%%%%%%%%%%%%%%%%%%%%
\definecolor{darkGreen}{rgb}{0,0.6,0}
\definecolor{gray}{rgb}{0.5,0.5,0.5}
\definecolor{mauve}{rgb}{0.58,0,0.82}
\definecolor{brightGreen}{rgb}{0,1,0}

%%%%%%%%%%%%%%%%%%%%%%%%%%%%%%%%%%%%%%%%%%%%%%%%%%%%%%%%%%%%%%%%%%%%%%%%%%%%%%%%
% Source code formatting
%%%%%%%%%%%%%%%%%%%%%%%%%%%%%%%%%%%%%%%%%%%%%%%%%%%%%%%%%%%%%%%%%%%%%%%%%%%%%%%%
% Default settings
\lstset{
    basicstyle=\ttfamily,           % font
                                    %
    title=\lstname,                 % show the filename of files included with \lstinputlisting
    captionpos=b,                   % sets the caption-position to bottom
    backgroundcolor=\color{white},  % choose the background color
                                    %
    numbers=left,                   % where to put the line-numbers
    numberstyle=\tiny\color{gray},  % the style that is used for the line-numbers
    stepnumber=1,                   % the step between two line-numbers
    numbersep=5pt,                  % how far the line-numbers are from the code
                                    %
    showspaces=false,               % show spaces adding particular underscores
    showstringspaces=false,         % underline spaces within strings
    showtabs=false,                 % show tabs within strings adding particular underscores
    tabsize=4,                      % sets default tabsize
    breaklines=true,                % sets automatic line breaking
    breakatwhitespace=false,        % sets if automatic breaks should only happen at whitespace
                                    %
    keywordstyle=\color{blue},      % keyword style
    commentstyle=\color{darkGreen}, % comment style
    stringstyle=\color{mauve},      % string literal style
    escapeinside={\%*}{*)},         % if you want to add a comment within your code
    morekeywords={*,...}            % if you want to add more keywords to the set
}
\lstloadlanguages{C,Matlab}

%%%%%%%%%%%%%%%%%%%%%%%%%%%%%%%%%%%%%%%%%%%%%%%%%%%%%%%%%%%%%%%%%%%%%%%%%%%%%%%%
% Variables (to save retyping and reformatting these throughout the document).
%%%%%%%%%%%%%%%%%%%%%%%%%%%%%%%%%%%%%%%%%%%%%%%%%%%%%%%%%%%%%%%%%%%%%%%%%%%%%%%%
%%%%%%%%%%%%%%%%%%%%%%%%%%%%%%%%%%%%%%%%%%%%%%%%%%%%%%%%%%%%%%%%%%%%%%%%%%%%%%%%
% Variables (to save retyping and reformatting these throughout the document).
%%%%%%%%%%%%%%%%%%%%%%%%%%%%%%%%%%%%%%%%%%%%%%%%%%%%%%%%%%%%%%%%%%%%%%%%%%%%%%%%
\usepackage{definevariable} % for \defineVariable
\usepackage{persons} % for \definePerson, \getPerson, \formatPerson

\defineVariable\thesis{}{Thesis}
\defineVariable\thesisTitle{\textsc}{Implementation of Random Projection Algorithms Using an FPGA}
\defineVariable\thesisDate{}{2012}
\defineVariable\studentName{\formatPerson}{\getPerson*{Spence}}
\defineVariable\supervisorName{\formatPerson}{\getPerson*{Leong}}

\defineVariable\degreeName{}{Bachelor of Engineering (Computer)}
\defineVariable\schoolName{}{School of Electrical Engineering}
\defineVariable\studentID{}{308216350}
\defineVariable\universityName{}{The University of Sydney}
\defineVariable\uosAssessment{}{\thesis}
\defineVariable\uosCode{}{ELEC4712/ELEC4713}

\urldef{\gitRepoHTTP}\url{https://github.com/joshuaspence/ThesisCode.git}
\urldef{\gitRepoSSH}\url{git@github.com:joshuaspence/ThesisCode.git}

% People
\definePerson{Chawla}{Sanjay Chawla}
\definePerson{Frechtling}{Michael Frechtling}
\definePerson{Heinrichs}{Amelia Heinrichs}
\definePerson{Khoa}{Nguyen Lu Dang Khoa}
\definePerson{Leong}{Philip Leong}
\definePerson{Spence}{Joshua Spence}

% PDF metadata
\newcommand*\pdfmetadata{%
    pdfauthor={Joshua Spence},
    pdftitle={Implementation of Random Projection Algorithms Using an FPGA},
    pdfsubject={},
    pdfkeywords={commute time,anomaly detection,FPGA},
    pdfdisplaydoctitle={true}
}

% Copyrights and trademarks
\newcommand\IntelXeon{Intel\textsuperscript{\textregistered} Xeon\textsuperscript{\textregistered}}
\newcommand\IntelCoreiSeven{Intel\textsuperscript{\textregistered} Core\textsuperscript{\texttrademark} i7}
\newcommand\AutoESL{\software{Xilinx\textsuperscript{\textregistered} AutoESL}}
\newcommand\ModelSim{\software{Altera\textsuperscript{\textregistered} ModelSim}}
\newcommand\Xilinx{Xilinx\textsuperscript{\textregistered}}
\newcommand\ARM{ARM\textsuperscript{\textregistered}}
\newcommand\Cortex{Cortex\textsuperscript{\texttrademark}}

%%%%%%%%%%%%%%%%%%%%%%%%%%%%%%%%%%%%%%%%%%%%%%%%%%%%%%%%%%%%%%%%%%%%%%%%%%%%%%%%
% Title page options
%%%%%%%%%%%%%%%%%%%%%%%%%%%%%%%%%%%%%%%%%%%%%%%%%%%%%%%%%%%%%%%%%%%%%%%%%%%%%%%%
\ifcsname beamer@author\endcsname
    \title[Random Projection Algorithms Using an FPGA]{Implementation of Random Projection Algorithms Using an FPGA}
    \author[Joshua Spence]{Joshua Spence\\{\tiny{Supervisor: Philip Leong}}}
    \institute[The University of Sydney]{
        School of Electrical Engineering\\
        The University of Sydney
    }
    \titlegraphic{\includegraphics[width=0.25\textwidth]{sydney_uni_coat_of_arms}}
\else
    \title{Implementation of Random Projection Algorithms Using an FPGA}
    \author{Joshua Spence}
\fi

\date{2012}

%%%%%%%%%%%%%%%%%%%%%%%%%%%%%%%%%%%%%%%%%%%%%%%%%%%%%%%%%%%%%%%%%%%%%%%%%%%%%%%%
% Custom commands
%%%%%%%%%%%%%%%%%%%%%%%%%%%%%%%%%%%%%%%%%%%%%%%%%%%%%%%%%%%%%%%%%%%%%%%%%%%%%%%%

%===============================================================================
% Format a block size profiling plot.
%-------------------------------------------------------------------------------
% \blockSizeProfilingPlot{DATASET}
%===============================================================================
\makeatletter
\newcommand\blockSizeProfilingPlot[1]{%
    \def\profilingPlots@width{0.9\textwidth}%
    \def\profilingPlots@height{0.9\textheight}%
    %
    \makebox[\textwidth][c]{%
        \maxsizebox{\profilingPlots@width}{\profilingPlots@height}{%
            \input{plots/block_size/#1}%
        }%
    }%
}
\makeatother

%===============================================================================
% Format a profiling plot.
%-------------------------------------------------------------------------------
% \profilingPlots{DATASET1}{DATASET2}{DATASET3}
%===============================================================================
\makeatletter
\newcommand\profilingPlots[3]{%
    \def\profilingPlots@verticalOffset{\vskip-1.5em}%
    \def\profilingPlots@columnWidth{0.33\textwidth}%
    \def\profilingPlots@pcaWidth{0.9\textwidth}%
    \def\profilingPlots@plotWidth{0.7\textwidth}%
    \def\profilingPlots@legendHeight{0.15\textheight}%
    %
    \profilingPlots@verticalOffset
    \begin{columns}[c]%
        \column{\profilingPlots@columnWidth}%
        \ifthenelse{\equal{#1}{}}{}{
            \begin{figure}%
                \includegraphics[width=\profilingPlots@pcaWidth]{pca/#1}\\%
                \resizebox{\profilingPlots@plotWidth}{!}{\input{plots/c/#1}}%
                \caption{\escape{#1}}%
            \end{figure}%
        }%
        %
        \column{\profilingPlots@columnWidth}%
        \ifthenelse{\equal{#2}{}}{}{
            \begin{figure}%
                \includegraphics[width=\profilingPlots@pcaWidth]{pca/#2}\\%
                \resizebox{\profilingPlots@plotWidth}{!}{\input{plots/c/#2}}%
                \caption{\escape{#2}}%
            \end{figure}%
        }%
        %
        \column{\profilingPlots@columnWidth}%
        \ifthenelse{\equal{#3}{}}{}{
            \begin{figure}%
                \includegraphics[width=\profilingPlots@pcaWidth]{pca/#3}\\%
                \resizebox{\profilingPlots@plotWidth}{!}{\input{plots/c/#3}}%
                \caption{\escape{#3}}%
            \end{figure}%
        }%
    \end{columns}%
    \vfill
    \makebox[\textwidth][c]{%
        \resizebox{!}{\profilingPlots@legendHeight}{%
            \pielegend[bound,color={yellow!60,blue!60,white!60,orange!60}]{
1/\escape{top_n_outlier_pruning_block},
1/\escape{distance_squared},
1/\escape{Other},
1/\escape{add_neighbour}
}
%
        }%
    }%
}
\makeatother

%%%%%%%%%%%%%%%%%%%%%%%%%%%%%%%%%%%%%%%%%%%%%%%%%%%%%%%%%%%%%%%%%%%%%%%%%%%%%%%%
% Custom formatting commands
%%%%%%%%%%%%%%%%%%%%%%%%%%%%%%%%%%%%%%%%%%%%%%%%%%%%%%%%%%%%%%%%%%%%%%%%%%%%%%%%

%===============================================================================
% Format the name of an algorithm.
%-------------------------------------------------------------------------------
% \algm{}
%===============================================================================
\newcommand\algm[1]{\emph{#1}}

%===============================================================================
% Format the name of a command line command. Also allows words to be broken over
% multiple lines
%
% See: http://jevopi.blogspot.com.au/2010/04/latex-lstinline-and-hyphenations.html
%-------------------------------------------------------------------------------
% \command{CMD}
%===============================================================================
\newcommand\command[1]{\lstinline[breaklines,literate={\_}{}{0\discretionary{\_}{}{\_}}]`#1`}

%===============================================================================
% Format the name of a data set.
%-------------------------------------------------------------------------------
% \dataset{DATASET}
%===============================================================================
\newcommand\dataset[1]{`\escape{#1}'}

%===============================================================================
% Format a computer host name.
%-------------------------------------------------------------------------------
% \hostname{HOSTNAME}
%===============================================================================
\newcommand\hostname[1]{\emph{#1}}

%===============================================================================
% Format the name of a programming language.
%-------------------------------------------------------------------------------
% \progLang{LANG}
%===============================================================================
\newcommand\progLang[1]{\lstinline`#1`}

%===============================================================================
% Format the name of a software package.
%-------------------------------------------------------------------------------
% \software{NAME}[CMD]
%===============================================================================
\NewDocumentCommand\software{om}{``#2''\IfNoValueTF{#1}{}{\space(\command{#1})}}

%%%%%%%%%%%%%%%%%%%%%%%%%%%%%%%%%%%%%%%%%%%%%%%%%%%%%%%%%%%%%%%%%%%%%%%%%%%%%%%%
% Graphics directory
%%%%%%%%%%%%%%%%%%%%%%%%%%%%%%%%%%%%%%%%%%%%%%%%%%%%%%%%%%%%%%%%%%%%%%%%%%%%%%%%
\graphicspath{{../img/}}

%%%%%%%%%%%%%%%%%%%%%%%%%%%%%%%%%%%%%%%%%%%%%%%%%%%%%%%%%%%%%%%%%%%%%%%%%%%%%%%%
% Bibliography
%%%%%%%%%%%%%%%%%%%%%%%%%%%%%%%%%%%%%%%%%%%%%%%%%%%%%%%%%%%%%%%%%%%%%%%%%%%%%%%%
\usepackage[backend=bibtex,language=australian]{biblatex} % use biblatex instead of bibtex
\bibliography{thesis}

%%%%%%%%%%%%%%%%%%%%%%%%%%%%%%%%%%%%%%%%%%%%%%%%%%%%%%%%%%%%%%%%%%%%%%%%%%%%%%%%
%                               END OF PREAMBLE                                %
%%%%%%%%%%%%%%%%%%%%%%%%%%%%%%%%%%%%%%%%%%%%%%%%%%%%%%%%%%%%%%%%%%%%%%%%%%%%%%%%
\begin{document}

%%%%%%%%%%%%%%%%%%%%%%%%%%%%%%%%%%%%%%%%%%%%%%%%%%%%%%%%%%%%%%%%%%%%%%%%%%%%%%%%
% Title Page
%%%%%%%%%%%%%%%%%%%%%%%%%%%%%%%%%%%%%%%%%%%%%%%%%%%%%%%%%%%%%%%%%%%%%%%%%%%%%%%%
\frame[label=titlepage,shrink]{\titlepage}

%%%%%%%%%%%%%%%%%%%%%%%%%%%%%%%%%%%%%%%%%%%%%%%%%%%%%%%%%%%%%%%%%%%%%%%%%%%%%%%%
% Introduction
%%%%%%%%%%%%%%%%%%%%%%%%%%%%%%%%%%%%%%%%%%%%%%%%%%%%%%%%%%%%%%%%%%%%%%%%%%%%%%%%
\section{Introduction}
In this chapter, I recount the design process that was adopted in order to
transform the anomaly detection using commute time algorithm into a hardware
design. Specifically, this chapter analyses software implementations of the
algorithm (in both \programmingLanguage{MATLAB} and \programmingLanguage{C}), in
order to identify performance bottlenecks that the algorithm faces on a typical
microprocess-based system.

Essentially, this chapter details a range of tests that were performed at
various stages along the development process in order to evaluate and compare
the performance of several variations of the basic algorithm implementation.

%%%%%%%%%%%%%%%%%%%%%%%%%%%%%%%%%%%%%%%%%%%%%%%%%%%%%%%%%%%%%%%%%%%%%%%%%%%%%%%%
% Base Implementation and Design Flow
%%%%%%%%%%%%%%%%%%%%%%%%%%%%%%%%%%%%%%%%%%%%%%%%%%%%%%%%%%%%%%%%%%%%%%%%%%%%%%%%
\subsection{Base Implementation and Design Flow}
\label{software:baseImplementation}
The base implementation of the anomaly detection using commute time algorithm,
on which much of the research of this Thesis is based, was generously provided
by \citeauthor{Khoa:2012}. \citeauthor{Khoa:2012} had designed a
\programmingLanguage{MATLAB} implementation of the algorithm as a part of his
research for \fullcite{Khoa:2012} \cite{Khoa:2012}.

This \programmingLanguage{MATLAB} code formed the foundation on which all
successive code of this Thesis eventuates. The design flow that was adopted for
the development process is as follows:
\begin{enumerate}
    \item \programmingLanguage{MATLAB} $\longrightarrow$
        \programmingLanguage{MATLAB}+\programmingLanguage{C} (using
        \programmingLanguage{MEX} files).
    \item \programmingLanguage{MATLAB}+\programmingLanguage{C} $\longrightarrow$
        Pure \programmingLanguage{C}.
    \item Pure \programmingLanguage{C} $\longrightarrow$ Pure
        \programmingLanguage{C} with \software{AutoESL} directives for \gls{HDL}
        synthesis.
    \item  Pure \programmingLanguage{C} with \software{AutoESL} directives
        $\longrightarrow$ pcore.
    % TODO
\end{enumerate}

% TODO: Flow chart

%%%%%%%%%%%%%%%%%%%%%%%%%%%%%%%%%%%%%%%%%%%%%%%%%%%%%%%%%%%%%%%%%%%%%%%%%%%%%%%%
% Data Sets Used For Testing and Profiling
%%%%%%%%%%%%%%%%%%%%%%%%%%%%%%%%%%%%%%%%%%%%%%%%%%%%%%%%%%%%%%%%%%%%%%%%%%%%%%%%
\subsection{Data Sets Used For Testing and Profiling}
\label{software:datasets}

%%%%%%%%%%%%%%%%%%%%%%%%%%%%%%%%%%%%%%%%%%%%%%%%%%%%%%%%%%%%%%%%%%%%%%%%%%%%%%%%
% The Algorithm
%%%%%%%%%%%%%%%%%%%%%%%%%%%%%%%%%%%%%%%%%%%%%%%%%%%%%%%%%%%%%%%%%%%%%%%%%%%%%%%%
\section{The Algorithm}
%%%%%%%%%%%%%%%%%%%%%%%%%%%%%%%%%%%%%%%%%%%%%%%%%%%%%%%%%%%%%%%%%%%%%%%%%%%%%%%%
% Introduction
%%%%%%%%%%%%%%%%%%%%%%%%%%%%%%%%%%%%%%%%%%%%%%%%%%%%%%%%%%%%%%%%%%%%%%%%%%%%%%%%
\subsection{Introduction}
\begin{frame}<1-3>{Introduction}{Anomaly Detection Using Commute Time}
    {\Large \textbf{Anomaly Detection Using Commute Time}}\\

    \medskip
    {\small \fullcite{Khoa:2012}}
    \note<1>{The algorithm that was selected for implementation was developed by
        ``Nguyen Lu Dang Khoa'' and published in ``Large Scale Anomaly Detection
        and Clustering Using Random Walks''.}

    \medskip
    \begin{columns}[c]
        \column{0.5\textwidth}
        \begin{figure}[H]
            \includegraphics[width=0.7\textwidth]{algorithm/run-time}
            \caption{Execution time of approximate commute time method versus
                Euclidean distance and local outlier factor methods}
        \end{figure}

        \column{0.5\textwidth}
        \begin{figure}[H]
            \includegraphics[width=0.7\textwidth]{algorithm/accuracy}
            \caption{Accuracy and execution time of approximated commute
                distance method compared to non-approximated method}
        \end{figure}
    \end{columns}

    \note<1>{These graphs show the measured performance of the chosen
        algorithm.}
    \note<2>{The left graph shows how run time (using \emph{approximated}
        commute time) is proportional to $O(n \log n)$, compared to:
        \begin{itemize}
            \item $O(n)$ using Euclidean distance
            \item $O(n^2)$ using local outlier factor
        \end{itemize}}
    \note<3>{The right graph shows how anomaly detection using
        \emph{approximated} commute time preserves a high percentage (84.6\% on
        average) of the top anomalies discovered without using approximations.}
\end{frame}

%%%%%%%%%%%%%%%%%%%%%%%%%%%%%%%%%%%%%%%%%%%%%%%%%%%%%%%%%%%%%%%%%%%%%%%%%%%%%%%%
% The Algorithm
%%%%%%%%%%%%%%%%%%%%%%%%%%%%%%%%%%%%%%%%%%%%%%%%%%%%%%%%%%%%%%%%%%%%%%%%%%%%%%%%
\subsection{The Algorithm}
\begin{frame}<1-2>[label=animation]{Animation}{A short animation illustrating the algorithm execution}
    \centering
    \movie[once,showcontrols]
        {\includegraphics[width=0.8\textwidth]{pca/testoutrank}}
        {animation/testoutrank.avi}

    \note<1>{This slide shows a short animation, illustrating the execution of
        the chosen anomaly detection method on a data set of 441 2-dimensional
        vectors.

        In this animation, the yellow points indicate the current block that is
        being processed. Each vector within the block is compared to all other
        vectors in the data set, illustrated by the blue marker in the
        animation.

        Each vector in the current block is given a score, a monotonic
        decreasing function that provides a heuristic to estimate if any given
        vector may be an outlier.}

    \note<2>{The current `best' outliers are highlighted in green in the
        animation. As the algorithm processes further blocks of data, the set of
        current outliers converges to the final set of outliers.

        If the score for any vector in the current block becomes less than the
        lowest score from the set of current outliers, then this vector is no
        longer considered to be a possible outlier, and is pruned from the data
        set. The pruned vectors appear in red in the animation.}
\end{frame}


%%%%%%%%%%%%%%%%%%%%%%%%%%%%%%%%%%%%%%%%%%%%%%%%%%%%%%%%%%%%%%%%%%%%%%%%%%%%%%%%
% Software Implementation
%%%%%%%%%%%%%%%%%%%%%%%%%%%%%%%%%%%%%%%%%%%%%%%%%%%%%%%%%%%%%%%%%%%%%%%%%%%%%%%%
\section{Software Implementation}
%%%%%%%%%%%%%%%%%%%%%%%%%%%%%%%%%%%%%%%%%%%%%%%%%%%%%%%%%%%%%%%%%%%%%%%%%%%%%%%%
% Algorithm Profiling
%%%%%%%%%%%%%%%%%%%%%%%%%%%%%%%%%%%%%%%%%%%%%%%%%%%%%%%%%%%%%%%%%%%%%%%%%%%%%%%%
\subsection{Algorithm Profiling}
\begin{frame}{Algorithm Profiling}{\escape{TopN_Outlier_Pruning_Block}}
    \profilingPlots{testoutrank}{runningex40k}{connect4}
\end{frame}
\begin{frame}{Algorithm Profiling}{\escape{TopN_Outlier_Pruning_Block}}
    \profilingPlots{testCDST}{testCDST2}{testCDST3}
\end{frame}
\begin{frame}{Algorithm Profiling}{\escape{TopN_Outlier_Pruning_Block}}
    \profilingPlots{spam_train}{spam}{mesh_network}
\end{frame}
\begin{frame}{Algorithm Profiling}{\escape{TopN_Outlier_Pruning_Block}}
    \profilingPlots{segmentation}{magicgamma}{connect4}
\end{frame}
\begin{frame}{Algorithm Profiling}{\escape{TopN_Outlier_Pruning_Block}}
    \profilingPlots{letter-recognition}{pendigits}{musk}
\end{frame}
\begin{frame}{Algorithm Profiling}{\escape{TopN_Outlier_Pruning_Block}}
    \profilingPlots{testoutrank}{testCD}{ball1}
\end{frame}
\begin{frame}{Algorithm Profiling}{\escape{TopN_Outlier_Pruning_Block}}
    \profilingPlots{runningex1k}{runningex10k}{runningex20k}
\end{frame}
\begin{frame}{Algorithm Profiling}{\escape{TopN_Outlier_Pruning_Block}}
    \profilingPlots{runningex30k}{runningex40k}{} % TODO: runningex50k
\end{frame}

%%%%%%%%%%%%%%%%%%%%%%%%%%%%%%%%%%%%%%%%%%%%%%%%%%%%%%%%%%%%%%%%%%%%%%%%%%%%%%%%
% Block Size Profiling
%%%%%%%%%%%%%%%%%%%%%%%%%%%%%%%%%%%%%%%%%%%%%%%%%%%%%%%%%%%%%%%%%%%%%%%%%%%%%%%%
\subsection{Block Size Profiling}
\begin{frame}[label=block-size-legend]{Block Size Profiling Plots}{Legend}
    \blockSizeProfilingPlot{legend}
\end{frame}

\begin{frame}{Total Execution Time}{\escape{commute_distance_anomaly}}
    \blockSizeProfilingPlot{total_execution_time}
\end{frame}

\begin{frame}{Total Run Time Complexity}{\escape{commute_distance_anomaly}}
    \begin{columns}[c]
        \column{0.5\textwidth}
        \blockSizeProfilingPlot{total_run_time_complexity.lin}

        \column{0.5\textwidth}
        \blockSizeProfilingPlot{total_run_time_complexity.log}
    \end{columns}
\end{frame}

\begin{frame}{Distance Calls}{\escape{TopN_Outlier_Pruning_Block}}
    \blockSizeProfilingPlot{distance_calls}
\end{frame}

\begin{frame}{Vectors Pruned}{\escape{TopN_Outlier_Pruning_Block}}
    \blockSizeProfilingPlot{vectors_pruned}
\end{frame}


%%%%%%%%%%%%%%%%%%%%%%%%%%%%%%%%%%%%%%%%%%%%%%%%%%%%%%%%%%%%%%%%%%%%%%%%%%%%%%%%
% Hardware Implementation
%%%%%%%%%%%%%%%%%%%%%%%%%%%%%%%%%%%%%%%%%%%%%%%%%%%%%%%%%%%%%%%%%%%%%%%%%%%%%%%%
\section{Hardware Implementation}
%%%%%%%%%%%%%%%%%%%%%%%%%%%%%%%%%%%%%%%%%%%%%%%%%%%%%%%%%%%%%%%%%%%%%%%%%%%%%%%%
% Design
%%%%%%%%%%%%%%%%%%%%%%%%%%%%%%%%%%%%%%%%%%%%%%%%%%%%%%%%%%%%%%%%%%%%%%%%%%%%%%%%
% In this section, illustrate the hardware and software interface for the
% design. Note that the most important part of the hardware design is
% parallelism, so it is important to mention how this was achieved.
\subsection{Design}
\begin{frame}[label=design]{Features}
    \begin{itemize}
        \item Majority of \escape{TopN_Outlier_Pruning_Block} function performed
            in software.

        \item Computation of expensive \escape{distance_squared} function
            performed in hardware.
        \begin{itemize}
            \item Multiple functional units.
            \item Pipelined design.
            \item Separate hardware units for loading ``outer'' and ``inner''
                vectors through DMA bus.
        \end{itemize}

        \item Software algorithm modified so as to utilise hardware
            implementation.
        \begin{itemize}
            \item Software can begin to process the next block element whilst
                waiting for \escape{distance_squared} computation to complete.
            \item Thread management and interrupt facilities required (with
                additional overhead).
            \item Increased instruction-level parallelism.
        \end{itemize}
    \end{itemize}
\end{frame}

\begin{frame}[label=design-diagram]{Diagram}
    \maxsizebox{\textwidth}{\textheight}{\begin{tikzpicture}[
            >=latex,
            node distance = 4em,
            interconnect/.style = {
                rectangle,
                draw,
                text centered,
                rounded corners,
                minimum width = 60em,
                minimum height = 5em
                },
            pcore/.style = {
                rectangle,
                draw,
                text centered,
                rounded corners,
                text width = 10em,
                minimum height = 2em,
                },
            dma/.style = {
                rectangle,
                draw,
                text centered,
                rounded corners,
                text width = 10em,
                minimum height = 2em,
                },
            arrow/.style = {
                draw,
                -latex'},
            auto]
         \node [interconnect] (axi) {AXI Interconnect Block};

         \node [dma, below=of axi.east, anchor=east] (dma-inner1) {AXI DMA};
         \node [pcore,left=of dma-inner1] (pcore-inner1) {AutoESL AXI Lite PCORE};
         \path [arrow] (dma-inner1.west) -| (pcore-inner1.east);

         \node [dma, below=of dma-inner1.east, anchor=east] (dma-inner2) {AXI DMA};
         \node [pcore,left=of dma-inner2] (pcore-inner2) {AutoESL AXI Lite PCORE};

         \node [dma, below=of dma-inner2.east, anchor=east] (dma-inner3) {AXI DMA};
         \node [pcore,left=of dma-inner3] (pcore-inner3) {AutoESL AXI Lite PCORE};

         \node [pcore,below=of axi.west, anchor=west] (pcore-outer) {AutoESL AXI Lite PCORE};,
         \node [dma, right=of pcore-outer] (dma-outer) {AXI DMA};
         \draw[thick,dotted] ($(pcore-outer.north west)+(-0.5,0.5)$) rectangle ($(dma-outer.south east)+(0.5,-0.5)$);
    \end{tikzpicture}}
\end{frame}

%%%%%%%%%%%%%%%%%%%%%%%%%%%%%%%%%%%%%%%%%%%%%%%%%%%%%%%%%%%%%%%%%%%%%%%%%%%%%%%%
% Results
%%%%%%%%%%%%%%%%%%%%%%%%%%%%%%%%%%%%%%%%%%%%%%%%%%%%%%%%%%%%%%%%%%%%%%%%%%%%%%%%
\subsection{Results}
\begin{frame}[label=performance-considerations]{Performance Considerations}
    \begin{itemize}
        \item Cannot transfer data faster than $O(n)$
        \item For each vector that is pruned from the block, we save (on
            average) $\frac{block\_size}{2}$ distance computations
    \end{itemize}
\end{frame}

\begin{frame}[label=performance-estimates]{Performance Estimates}
    The execution time for the proposed design is estimated to be the sum of:
    \begin{enumerate}[<+->]
        \item Execution time of software implementation, excluding the
            distance\_squared function
        \item Time required for DMA transfer of ``outer'' vector.
        \item Time required for DMA transfer of ``inner'' vector, plus the time
            required for the hardware to complete the computation and return the
            result
        \item Overheads associated with thread management
    \end{enumerate}
\end{frame}

%%%%%%%%%%%%%%%%%%%%%%%%%%%%%%%%%%%%%%%%%%%%%%%%%%%%%%%%%%%%%%%%%%%%%%%%%%%%%%%%
% Discussion
%%%%%%%%%%%%%%%%%%%%%%%%%%%%%%%%%%%%%%%%%%%%%%%%%%%%%%%%%%%%%%%%%%%%%%%%%%%%%%%%
\subsection{Discussion}
\begin{frame}[label=discussion]{Discussion}
    \begin{itemize}
        \item Hardware design takes X cycles

        \item On this hardware, we expect a speed-up of X\%

        \item Maximum expected improvement (using Amdahl's Law) is X\%
        \note[item]{Amdahl's Law is used to find the maximum expected
            improvement to an overall system when only part of the system is
            improved. It is often used in parallel computing to predict the
            theoretical maximum speed-up using multiple processors.}

        \item Scaling
    \end{itemize}
\end{frame}

%%%%%%%%%%%%%%%%%%%%%%%%%%%%%%%%%%%%%%%%%%%%%%%%%%%%%%%%%%%%%%%%%%%%%%%%%%%%%%%%
% Conclusion
%%%%%%%%%%%%%%%%%%%%%%%%%%%%%%%%%%%%%%%%%%%%%%%%%%%%%%%%%%%%%%%%%%%%%%%%%%%%%%%%
\section{Conclusion}
%%%%%%%%%%%%%%%%%%%%%%%%%%%%%%%%%%%%%%%%%%%%%%%%%%%%%%%%%%%%%%%%%%%%%%%%%%%%%%%%
% The Function
%%%%%%%%%%%%%%%%%%%%%%%%%%%%%%%%%%%%%%%%%%%%%%%%%%%%%%%%%%%%%%%%%%%%%%%%%%%%%%%%
\section{The Function}
\label{profiling:function}
The function that was identified in the aforementioned tests as being most
significant to the execution time of the \algm{anomaly detection using commute
time} algorithm, and thus an ideal candidate for hardware acceleration, was the
\command{Top_N_Outlier_Pruning_Block} function. This function was originally
proposed by \citeauthor{Bay:2003} in \citetitle{Bay:2003}.

Essentially, this function simply performs a nested loops comparison of the
distances between pairwise vectors. Consequently, in the worst case, the
performance of this algorithm is very poor --- and it could require $O(n^2)$
distance computations and $O(\frac{n}{blocksize} \times n)$ data accesses
\cite{Bay:2003}. However, the use of randomization as well as pruning improves
the average-case complexity of the algorithm to be somewhat near-linear.

The main idea of the nested loop algorithm is that for each vector in the data
set, we keep track of the closest neighbours found so far. When an vector's
closest neighbours achieve a score lower than the cutoff, the data set is pruned
from the current block because it can no longer be an outlier. As the algorithm
proceeds to process more vectors, it is able to find more extreme outliers and
the cutoff increases along with pruning efficiency \cite{Bay:2003}. The score
function can be any monotonically decreasing function of the nearest neighbour
distances such as the distance to the $k$th nearest neighbour, or the average
distance to the k neighbours \cite{Bay:2003}.

In \citetitle{Bay:2003}, \citeauthor{Bay:2003} proves that the performance of
the \command{Top_N_Outlier_Pruning_Block} function was capable of scaling to
large, real, high-dimensional data sets, whilst maintaining a near linear
scaling performance. However, the algorithm depends on a number of assumptions,
violations of which can lead to poor performance.

\begin{enumerate}
    \item The algorithm assumes that the data is in random order. If the data is
        not in random order and is sorted then the performance can be poor.
    \item The outlier algorithm depends on the independence of examples. If
        examples are dependent in such a way that they have similar values (and
        will likely be in the set of $k$ nearest neighbours) this can cause
        performance to be poor as the algorithm may have to scan the entire data
        set to find the dependent examples.
    \item The third situation when our algorithm can perform poorly occurs when
        the data does not contain outliers. For example, our experiment with the
        examples drawn from a uniform distribution had very poor scaling.
        However, we believe data sets of this type are likely to be rare as most
        physical quantities one can measure have distributions with tails.
\end{enumerate}

\begin{algorithm}
    \LinesNumbered

\SetKwInput{InputK}{k}\InputK{the number of nearest neighbours to consider}
\SetKwInput{InputN}{N}\InputN{the number of outliers to return}
\SetKwInput{InputData}{data}\InputData{a set of examples in random order}
\SetKwInOut{Output}{outliers}\Output{a set of outliers}
\SetKwInOut{OutputScores}{outlier\_scores}\OutputScores{}

\SetKwData{BlockIterator}{b}
\SetKwData{Block}{block}
\SetKwData{Cutoff}{cutoff}
\SetKwData{DataIterator}{d}
\SetKwData{Data}{data}
\SetKwData{NumK}{k}
\SetKwData{NumN}{N}
\SetKwData{Neighbours}{neighbours}
\SetKwData{OutlierIterator}{o}
\SetKwData{Outliers}{outliers}
\SetKwData{Score}{score}

\SetKwFunction{Closest}{closest}
\SetKwFunction{Distance}{distance}
\SetKwFunction{LoadBlock}{load\_next\_block}
\SetKwFunction{MaxDist}{max\_dist}
\SetKwFunction{Min}{min}
\SetKwFunction{Top}{top}

\Begin{
    \Cutoff $\leftarrow$ 0\tcp*[l]{set the cutoff for pruning to $0$}
    \Outliers $\leftarrow\emptyset$\tcp*[l]{initialize to the empty set}
    \BlankLine
    \While(\tcp*[h]{load a block of examples from \Data}){\Block $\leftarrow$ \LoadBlock{\Data}}{
        \ForEach{\BlockIterator $\in$ \Block}{
            $\Neighbours[\BlockIterator]\leftarrow\emptyset$\;
        }
        \BlankLine
        \ForEach{\DataIterator $\in$ \Data}{
            \ForEach{\BlockIterator $\in$ \Block $:$ \BlockIterator $\neq$ \DataIterator}{
                \If{$|\Neighbours[\BlockIterator]|<$ \NumK $\:\lor\:$ \Distance{\BlockIterator, \DataIterator} $<$ \MaxDist{\BlockIterator, $\Neighbours[\BlockIterator]$}}{
                    %$\Neighbours[\BlockIterator] \leftarrow$ \Closest{\BlockIterator, $\Neighbours[\BlockIterator]$} $\cup$ \DataIterator,\NumK)$\;
                    \If{$\Score[\Neighbours[\BlockIterator],\BlockIterator)<\Cutoff$}{
                        \Block $\leftarrow$ \Block $\setminus$ \BlockIterator\;
                    }
                }
            }
        }
        \BlankLine
        \Outliers $\leftarrow$ \Top{\Block $\cup$ \Outliers, \NumN}\;
        \Cutoff $\leftarrow$ $\Min_{\OutlierIterator \in \Outliers}(\Score(\OutlierIterator))$\tcp*[l]{the cutoff is the score of the weakest outlier}
    }

    \BlankLine
    \KwRet{\Outliers}\;
}

    \caption{TopN\_Outlier\_Pruning\_Block}
    \label{algm:TopNOutlierPruningBlock}
\end{algorithm}

The \command{Top_N_Outlier_Pruning_Block} function is listed in
\autoref{algm:TopNOutlierPruningBlock}. It is worthwhile to note that the
\command{score} function can be any monotonically decreasing function of the
nearest neighbour distances such as the distance to the $k$th nearest neighbour,
or the average distance to the $k$ neighbours \cite{Bay:2003}.

The main idea of the \command{Top_N_Outlier_Pruning_Block} algorithm is that for
each example in the input data set, the algorithm keeps track of the closest
neighbours found so far. When an example's closest neighbours achieve a score
lower than the `cufoff', the example is discarded because it can no longer be an
outlier. As more examples are processed, the algorithm finds more extreme
outliers and the cutoff increases along with pruning efficiency \cite{Bay:2003}.

In the worst case, the performance of the algorithm is very poor --- due to the
nested loops, it could require $O(n^2)$ distance computations and
$O(\frac{n}{blocksize} \times n)$ data accesses. However, \citeauthor{Bay:2003}
proved (through application of the algorithm to both real and synthetic data
sets) that the algorithm performs considerably better than the expected
$O(n^2)$ running time in the average case. The performance improvements over
similar algorithms were attributed to the application of randomization and
pruning techniques. The outlier pruning problem can be considered similar to the
problem of conducting a set of independent Bernoulli trials in which examples
are analysed until $k$ examples within distance $d$ are found, or until the data
set is exhausted. \citeauthor{Bay:2003} proved that the number of trials
expected to achieve $k$ examples within distance $d$ is given by:
\begin{equation}
    E[Y] \leq \frac{k}{\pi(\textbf{x})} + \Bigg(1-\sum_{y=k}^{N} P(Y=y)\Bigg) \times N
    \label{eqn:outlierPruneComplexity}
\end{equation}
Where $\pi(x)$ is the probability that a random drawn example lies within
distance $d$ of point $x$ and $P(Y=y)$ is the probability of obtaining the $k$th
success on trial $y$.

The first term in \autoref{eqn:outlierPruneComplexity} represents the number of
distance computations to eliminate non-outliers, and does not depend on $n$. The
second term represents the expected cost of outliers, and does depend on $n$,
yielding an overall quadratic dependency to process $n$ examples in total.
However, note that we typically set the program parameters to return a small an
possibly fixed number of outliers. Thus the first term dominates and we obtain
near linear performance \cite{Bay:2003}. More specifically, it was determined
that the primary factor determining the scaling is how the cut-off changes as
$n$ increases.

%%%%%%%%%%%%%%%%%%%%%%%%%%%%%%%%%%%%%%%%%%%%%%%%%%%%%%%%%%%%%%%%%%%%%%%%%%%%%%%%
% Appendix
%%%%%%%%%%%%%%%%%%%%%%%%%%%%%%%%%%%%%%%%%%%%%%%%%%%%%%%%%%%%%%%%%%%%%%%%%%%%%%%%
\appendix
\chapter{Introduction}
\label{introduction}

%%%%%%%%%%%%%%%%%%%%%%%%%%%%%%%%%%%%%%%%%%%%%%%%%
% Motivation
%%%%%%%%%%%%%%%%%%%%%%%%%%%%%%%%%%%%%%%%%%%%%%%%%
\section{Motivation}
\label{motivation}
Anomaly detection is an important technique which can be applied to a wide range
of applications. There are many techniques to detect and measure anomalies, each
usually most applicable to some specific problem domain \citeNeeded{}. A general
algorithm for anomaly detection has proved difficult to design \citeNeeded{},
due to various challenges associated with defining and measuring anomalies. Most
existing approaches are based on statistical or geometrical measures involving
distance. Whilst many algorithms work well for some subset of input data sets,
it has been difficult to discover an algorithm which performs both correctly and
efficiently on all possible data sets.

Previous research has found merit in applying randomization techniques to highly
multivariate data sets in order to reduce the dimensionality of these data sets
whilst maintaining their fundamental and statistical properties. Such reduction
of the dimensionality of data, assuming it can be performed efficiently, allows
previously-unscalable anomaly detection algorithms to be practically applied to
a wider range of data and applications.

Anomaly detection is an important and contemporary problem in the field of
computer science, and is of particular interesting to stock market analysis,
network intrusion detection and image comparison.

%%%%%%%%%%%%%%%%%%%%%%%%%%%%%%%%%%%%%%%%%%%%%%%%%
% Contributions of this thesis
%%%%%%%%%%%%%%%%%%%%%%%%%%%%%%%%%%%%%%%%%%%%%%%%%
\section{Contributions of this thesis}
\label{contributions}
One key difficulty in anomaly detection is the efficient scaling of a general
algorithm to apply to highly multivariate data \citeNeeded{}. In this thesis, I
explore the use of randomization techniques (such as random projections and
commute time) in anomaly detection algorithms. These techniques provide
encouraging results with regards to the run-time complexity of an algorithm.

Furthermore, in this thesis I attempt to make observations and an analysis of
the run-time performance of anomaly detection algorithms using randomization
techniques, so as to identify steps in the algorithms that bottleneck the
algorithm's performance. Through this identification it would be possible to
improve the \emph{actual} run-time performance of the algorithm by utilising
the advantages of reconfigurable computing.


%%%%%%%%%%%%%%%%%%%%%%%%%%%%%%%%%%%%%%%%%%%%%%%%%
% Organization
%%%%%%%%%%%%%%%%%%%%%%%%%%%%%%%%%%%%%%%%%%%%%%%%%
\section{Organization}
\label{organization}
The rest of this \thesis{} is organized as follows. In \autoref{background}, I
provide a background to various anomaly detection techniques, as well as
randomization techniques. In order to provide the reader with an understanding
of the background topics, a brief background is given to various topics in
linear algebra, vector calculus and graph theory. However, the provided
background material is not intended to be comprehensive, and it is expected that
the reader is familiar with these fields of mathematics. In
\autoref{reconfigurableComputing}, I provide an overview of reconfigurable
computing, including an explanation of \glspl{FPGA} and an overview of
\gls{HLS}.

In \autoref{software}, I profile the execution of an anomaly detection algorithm
and explore possible improvements to the algorithm, in particular by outsourcing
various stages of the algorithm to an FPGA device. I proceed to make iterative
refinements to the algorithm from its base implementation towards a hardware
design.

In \autoref{hardware}, I propose several hardware implementations and discuss
design considerations, as well as the design process itself, that must be
considered for these designs. I also detail the results of some initial testing
investigating the performance expectations from differing hardware designs.

I conclude in \autoref{conclusions} by reflecting upon the results obtained
through this research, and making suggestions for further research in this
topic.

 
%%%%%%%%%%%%%%%%%%%%%%%%%%%%%%%%%%%%%%%%%%%%%%%%%
% Schedule
%%%%%%%%%%%%%%%%%%%%%%%%%%%%%%%%%%%%%%%%%%%%%%%%%
\section{Schedule}
\label{schedule}
A Gantt chart showing the anticipated schedule for the project is shown in 
\autoref{fig:ganttChart}. This will be updated as the project progresses.

% Gantt chart
% See http://www.martin-kumm.de/tex_gantt_package.php
\begin{figure}[h]
	\centering
	\scalebox{0.5}{
	\begin{gantt}[
			xunitlength=0.5cm,
			fontsize=\small,
			titlefontsize=\small,
			drawledgerline=true
			]{18}{36}
		% Months
		\begin{ganttitle}
			\titleelement{Mar}{4}
			\titleelement{Apr}{4}
			\titleelement{May}{4}
			\titleelement{Jun}{4}
			\titleelement{Jul}{4}
			\titleelement{Aug}{4}
			\titleelement{Sep}{4}
			\titleelement{Oct}{4}
			\titleelement{Nov}{4}
		\end{ganttitle}
		% Split months into quarters
		\begin{ganttitle}
			\numtitle{1}{1}{4}{1}
			\numtitle{1}{1}{4}{1}
			\numtitle{1}{1}{4}{1}
			\numtitle{1}{1}{4}{1}
			\numtitle{1}{1}{4}{1}
			\numtitle{1}{1}{4}{1}
			\numtitle{1}{1}{4}{1}
			\numtitle{1}{1}{4}{1}
			\numtitle{1}{1}{4}{1}
		\end{ganttitle}

		\input{ganttdata}
	\end{gantt}
	}
	\caption{Schedule for thesis work}
	\label{fig:ganttChart}
\end{figure}



\end{document}